% !TEX root = ../main.tex
\section{Derivations Supporting the Two-Stage Estimator}

We briefly restate the derivations that motivate the two-stage fitting procedure. Let $f(x; \bm{P})$ denote a complete-form probability density defined on $(0, \infty)$ with parameter vector $\bm{P}$. Observations are only available for the truncated interval $[a, b]$, with $a = 10$~cm and $b = 60$~cm in the PSP experiment. The truncated form of $f$ is
\begin{equation}
  f^{T}(x; \bm{P}) =
    \begin{cases}
      c \, f(x; \bm{P}), & a \leq x \leq b, \\
      0, & \text{otherwise},
    \end{cases}
    \label{eq:truncated-density}
\end{equation}
where the normalising constant is
\begin{equation}
  c = \left[\int_a^b f(t; \bm{P}) \,\mathrm{d}t\right]^{-1} = \left[F(b; \bm{P}) - F(a; \bm{P})\right]^{-1}.
\end{equation}
Here $F$ is the cumulative distribution function associated with $f$. When a closed-form expression for $F$ is unavailable, $c$ must be evaluated numerically, complicating routine model comparison.

To avoid deriving bespoke truncated forms for each target distribution, we augment the complete density with a free scaling parameter $s$ and define
\begin{equation}
  f'(x; \bm{P}', s) = s \, f(x; \bm{P}),
\end{equation}
where $\bm{P}' = \bm{P} \cup \{s\}$. Because $s$ relaxes the unity constraint on the integral of $f$, the least-squares objective
\begin{equation}
  Z(\bm{P}', s) = \sum_{i \in I} \left[f'(x_i; \bm{P}', s) - \hat{y}_i^{(n)}\right]^2
\end{equation}
admits solutions with residual behaviour comparable to $f^{T}$ over $[a, b]$. The augmented model reproduces any truncated shape achievable by $f^{T}$ but introduces an additional degree of freedom that inflates parameter variance.

If the truncated solution is matched exactly on $[a,b]$, then $f'(x)=f^{T}(x)$ implies $s=c=[F(b;\bm{P})-F(a;\bm{P})]^{-1}$. Hence $1/s$ is the implied retained mass on $[a,b]$, and departures from that value quantify any residual mass assigned outside the truncation interval.

We therefore fit the model in two stages. Stage~1 solves
\begin{equation}
  (\widehat{\bm{P}'}, \hat{s}) = \arg\min_{\bm{P}', s} Z(\bm{P}', s),
\end{equation}
yielding a consistent estimate $\hat{s}$ that scales the complete-form density to the truncated support. Stage~2 fixes the scaling parameter and re-optimises the original parameter vector,
\begin{equation}
  \widehat{\bm{P}} = \arg\min_{\bm{P}} Z(\bm{P}, \hat{s}),
\end{equation}
where $Z(\bm{P}, \hat{s})$ denotes the objective with $s$ held constant. The fitted curves coincide,
\begin{equation}
  Z(\widehat{\bm{P}'}, \hat{s}) \simeq Z(\widehat{\bm{P}}, \hat{s}),
\end{equation}
but the second-stage covariance matrix excludes the nuisance scaling parameter, yielding tighter standard errors for the distributional parameters of interest.

The approach generalises to any probability density for which numerical quadrature of $f$ over $[a, b]$ is feasible. In the manuscript we implement $f$ as the Weibull and gamma specialisations of the generalised gamma distribution,
\begin{equation}
  f_{\mathrm{GG}}(x; a, \beta, p) = \frac{a x^{ap - 1} \exp\left[-(x / \beta)^a\right]}{\beta^{ap} \Gamma(p)}, \qquad x > 0,
\end{equation}
with the Weibull recovered by fixing $p = 1$ and the gamma by fixing $a = 1$. The truncated form in \eqref{eq:truncated-density} follows by replacing $f$ with $f_{\mathrm{GG}}$ and evaluating the normalising constant $c$ numerically.
